\section{Proofs of results from Section~\ref{sect:main}}

We will repeatedly use the following straightforward technical fact about the root system of type~$A_{n-1}$, the proof
of which is left to the reader as an exercise.
\begin{claim}
\label{claim:cocycle}
Let $A$ be a~multiplicatively written abelian group.
Let $\phi_{ij}\colon \mathbb Z^n\to A$, $1\le i,j\le n$, $i\ne j$, be a~system of maps satisfying
\begin{gather*}
\phi_{ij}(w(\mathbf k)) = \phi_{w^{-1}(i),w^{-1}(j)}(\mathbf k),
\qquad
\phi_{ji}(\mathbf k)\phi_{ij}(\mathbf k)=1
\end{gather*}
for all $\mathbf k=(k_1,\ldots,k_n)\in\mathbb Z^n$, $w\in\mathbb S_n$, where $w(\mathbf k)$ stands for
$(k_{w^{-1}(1)},\ldots,k_{w^{-1}(n)})$.
Denote
\begin{gather*}
\phi_w(\mathbf k) = \prod\limits_{i<j,w(i)>w(j)} \phi_{ij}(\mathbf k).
\end{gather*}
Then:
\begin{enumerate}\itemsep=0pt
\item[$(a)$] For all $w',w\in \mathbb S_n$,
\begin{gather*}
\phi_{w'w}(\mathbf k) = \phi_{w'}(w(\mathbf k))\phi_w(\mathbf k).
\end{gather*}

\item[$(b)$] If for all $i\ne j$ and for all $\mathbf k,\mathbf{k'}\in \mathbb Z^n$ one has
\begin{gather*}
\phi_{ij}(\mathbf k+\mathbf{k'}) = a_{ij}^{k_ik_j'- k_i'k_j} \phi_{ij}(\mathbf k)\phi_{ij}(\mathbf{k'}),
\end{gather*}
then for all $w\in \mathbb S_n$ one has
\begin{gather*}
\langle \mathbf k,\mathbf{k'}\rangle \phi_{w}(\mathbf k+\mathbf{k'}) = \langle w(\mathbf k),w(\mathbf{k'})\rangle
\phi_{w}(\mathbf k)\phi_{w}(\mathbf{k'}),
\end{gather*}
where $\langle \mathbf k,\mathbf{k'}\rangle= \prod\limits_{i<j}a_{ij}^{k_i'k_j}$.
Here $a_{ij}$ are elements of $A$ such that $a_{ji}=a_{ij}$.
\end{enumerate}
\end{claim}

\subsection*{Proof of Proposition~\ref{prop:automorphism}}

Observe that the group $\mathbb G_n\subset \mathrm{GL}_n(\mathbb C)$ is generated by $\mathbb S_n$ and $(\mathbb
C^\times)^n$ subject to the semidirect product relations
\begin{gather*}
wt_j^{(\zeta)}=t_{w(j)}^{(\zeta)} w,
\qquad
\text{for all} \ \ w\in \mathbb S_n,
\quad
j\in\{1,\ldots,n\},
\quad
\zeta\in\mathbb C^\times.
\end{gather*}
In what follows, we use the abbreviation $x^{\mathbf k}$ to denote $x_1^{k_1}\cdots x_n^{k_n}$, where $\mathbf
k=(k_1,\ldots,k_n)\in \mathbb Z_{\ge 0}^n$.
For $c\in \mathbb C^\times$, def\/ine $\phi_{ij}^{(c)}\colon \mathbb Z^n\to \mathbb C^\times$ by
\begin{gather*}
\phi_{ij}^{(c)}(\mathbf k) = (-1)^{k_ik_j} c^{\overline k_i-\overline k_j},
\qquad
\text{where}
\qquad
\overline k=(k\ \mathrm{mod}\ 2)\in\{0,1\}.
\end{gather*}
Additionally, def\/ine $\phi_{ij}^{(0)}(\mathbf k)=1$ for all $\mathbf k$.
Because $\phi_{ij}^{(c^{-1})}(\mathbf k)=\phi_{ij}^{(c)}(\mathbf k)^{-1}=\phi_{ji}^{(c)}(\mathbf k)$, the system~$\phi_{ij}^{(c)}$ satisf\/ies the condition in Claim~\ref{claim:cocycle} and hence gives rise to functions
$\phi_w^{(c)}$.
The following lemma is then immediate from Claim~\ref{claim:cocycle}(a).

\begin{Lemma}
\label{lem:action}
For each $c\in\mathbb C$, the formula
\begin{gather*}
w t_1^{(\zeta_1)}\cdots t_n^{(\zeta_n)} \mathop{\triangleright}\nolimits_c x^{\mathbf k} = \phi^{(c)}_w(\mathbf k)
\left(\prod\limits_{j=1}^n \zeta_j^{k_j}\right) x^{w(\mathbf k)}
\end{gather*}
defines an action $\mathop{\triangleright}\nolimits_c$ of the group $\mathbb G_n$ on the space $S(V)$.
\end{Lemma}

The actions $\mathop{\triangleright}\nolimits_0$, $\mathop{\triangleright}\nolimits_1$ given by Lemma~\ref{lem:action}
coincide on the generators $s_i$, $t_j^{(\zeta)}$ of $\mathbb G_n$ with the actions
$\mathop{\triangleright}\nolimits_+$, $\mathop{\triangleright}\nolimits_-$ def\/ined in part~(a) of
Proposition~\ref{prop:automorphism}.
Hence part~(a) of the proposition is proved.



Denote by $\bullet_+$, respectively $\bullet_-$, the multiplication on the algebra $S(V)$, respectively
$S_{\mathbf{-1}}(V)$.
One has the following multiplication rule for monomials:
\begin{gather*}
x^{\mathbf k}\bullet_\pm x^{\mathbf{k'}} = \langle \mathbf k,\mathbf k'\rangle_\pm x^{\mathbf{k}+\mathbf{k'}}
\qquad
\text{for all}\quad\mathbf{k}, \mathbf{k'}\in \mathbb Z_{\ge 0}^n,
\end{gather*}
where $\langle \mathbf k,\mathbf k'\rangle_\pm$ is as given in Claim~\ref{claim:cocycle}(b) with all $a_{ij}=\pm1$, $i\ne j$.

It is enough to check that $w\mathop{\triangleright}\nolimits_\pm$ and
$t_j^{(\zeta)}\mathop{\triangleright}\nolimits_\pm$ are automorphisms of the respective algebra structures on $S(V)$.
We apply these actions to both sides of the multiplication rule for monomials and check that the results are equal.
This is trivial for $t_j^{(\zeta)}\mathop{\triangleright}\nolimits_\pm$.
For $w\mathop{\triangleright}\nolimits_\pm$ where $w\in \mathbb S_n$,   the equality is guaranteed by
Claim~\ref{claim:cocycle}(b) and Lemma~\ref{lem:action}, applied to functions $\phi_{ij}^+=\phi_{ij}^{(0)}$,
respectively $\phi_{ij}^-=\phi_{ij}^{(1)}$.
This proves parts (b), (c) of Proposition~\ref{prop:automorphism}.

Now let us prove part (d) of Proposition~\ref{prop:automorphism}.
Clearly, the natural $\mathbb S_n$-action on the group $(\mathbb C^\times)^ n\subset \mathbb G_n$ extends to that on the
group algebra $\mathbb C (\mathbb C^\times)^n$.

For each $c\in \mathbb C\setminus\{0\}$, $w\in \mathbb S_n$, def\/ine the element $Q_w^{(c)}\in \mathbb C (\mathbb C^\times)^n$ by
\begin{gather*}
Q_w^{(c)} = \prod\limits_{i<j:w(i)>w(j)} Q_{ij}^{(c)},
\end{gather*}
where $ Q_{ij}^{(c)} = \frac14 \big(\big(c+c^{-1}\big)\big(1-\g i\g j\big)+ \big(c-c^{-1}+2\big)\g i+\big(c^{-1}-c+2\big)\g j\big)$, $1\le i,j\le n$.
\begin{Lemma}\label{lem:Q}
For all $w',w\in \mathbb S_n$, $c\in\mathbb C^\times$, $Q_w^{(c)}Q_w^{(c^{-1})}=1$ and
$Q_{w'w}^{(c)}=w^{-1}(Q_{w'}^{(c)})\cdot Q_w^{(c)}$.
\end{Lemma}
\begin{proof}
Denote by $\mathcal F$ the algebra of all functions from $\mathbb Z^n$ to $\mathbb C$ with pointwise addition and
multiplication.
Clearly, the assignment $t_1^{(\zeta_1)}\ldots t_n^{(\zeta_n)} \mapsto (\mathbf k\mapsto \zeta_1^{k_1}\ldots
\zeta_n^{k_n})$ def\/ines a~group homomorphism $(\mathbb C^\times)^n\to \mathcal F^\times $ which extends to an injective
map $\Psi\colon \mathbb C(\mathbb C^\times)^n\hookrightarrow \mathcal F$.

It is easy to check that $\Psi(Q_{ij}^{(c)})=\phi_{ij}^{(c)}$ where $\phi_{ij}^{(c)}$ is as in Lemma~\ref{lem:action}.
In particular, $Q_{ij}^{(c)}Q_{ij}^{(c^{-1})}=1$ since $\phi_{ij}^{(c)}\phi_{ij}^{(c^{-1})}=1$ and $\Psi$ is injective.
This proves the f\/irst assertion of the lemma.
To prove the second assertion, apply $\Psi^{-1}$ to Claim~\ref{claim:cocycle}(a) and use the fact that the function
$\mathbf k\mapsto \phi(w(\mathbf k))$ is mapped by $\Psi^{-1}$ to $w^{-1}(\Psi^{-1}(\phi))$ for all functions $\phi$
from the subgroup of $\mathcal F^\times$ generated by $\{\phi_{ij}^{(c)}:i\ne j\}$.
\end{proof}

Now for each $c\in\mathbb C^\times$ def\/ine the $\mathbb C$-linear map $J_c\colon \mathbb C\mathbb G_n\to \mathbb
C\mathbb G_n$ by the formula
\begin{gather*}
J_c(wt)=wtQ^{(c)}_w,
\qquad
t\in \mathbb C(\mathbb C^\times)^n,
\qquad
w\in \mathbb S_n.
\end{gather*}
\begin{Lemma}\label{lem:J_c}
For each $n\ge 1$,
\begin{enumerate}\itemsep=0pt
\item[$(a)$] $J_c$ is an algebra automorphism of $\mathbb C \mathbb G_n$ with inverse $J_{c^{-1}}$.

\item[$(b)$] $\rho_+\circ J_c=\rho_c$, where $\rho_c:\mathbb C \mathbb G_n\to \End_\mathbb C S(V)$ is the algebra homomorphism
corresponding to $\mathop{\triangleright}\nolimits_c$.
\end{enumerate}
\end{Lemma}

\begin{proof}
On the one hand, $J_c(w't'w t)=J_c(w'w\cdot w^{-1}(t')t)=w'w \cdot w^{-1}(t')tQ_{w'w}^{(c)}$.
On the other hand,
\begin{gather*}
J_c(w't')J_c(w t) = w't'Q_{w'}^{(c)} wt Q^{(c)}_w=w'w \cdot w^{-1}(t')tw^{-1}(Q_{w'}^{(c)}) Q_w^{(c)}=J_c(w't'wt)
\end{gather*}
by the second assertion of Lemma~\ref{lem:Q}.
Hence $J_c$ is a~homomorphism of algebras.
Now{\samepage
\begin{gather*}
J_c(J_{c^{-1}}(wt))=J_c\big(wtQ_w^{(c^{-1})}\big)=wtQ_w^{(c^{-1})}Q_w^{(c)}=wt
\end{gather*}
by the f\/irst assertion of Lemma~\ref{lem:Q}.
This proves part (a) of the lemma.}

Prove (b).
In view of Lemma~\ref{lem:action}, it suf\/f\/ices to show that $Q_w^{(c)}\mathop{\triangleright}\nolimits_+ x^{\mathbf
k}=\phi_w^{(c)}(\mathbf k) x^{\mathbf k}$.
Indeed,
\begin{gather*}
Q_w^{(c)}\mathop{\triangleright}\nolimits_+ x^{\mathbf k}
= \frac14 \big(\big(c+c^{-1}\big)\big(1-(-1)^{k_i+k_j}\big)+\big(c-c^{-1}+2\big)(-1)^{k_i}+\big(c^{-1}-c+2\big)(-1)^{k_j}\big)x^{\mathbf k}
\\
\phantom{Q_w^{(c)}\mathop{\triangleright}\nolimits_+ x^{\mathbf k}}
=\phi_{ij}^{(c)}(\mathbf k)x^{\mathbf k}.
\end{gather*}

Finally,
\begin{gather*}
 J_c(wt)\mathop{\triangleright}\nolimits_+ x^{\mathbf k}
=\big(wtQ_w^{(c)}\big)\mathop{\triangleright}\nolimits_+x^{\mathbf k}
=(wt)\mathop{\triangleright}\nolimits_+\big(Q_w^{(c)}\mathop{\triangleright}\nolimits_+x^{\mathbf k}\big)
=\phi_w^{(c)}(\mathbf k)(wt)\mathop{\triangleright}\nolimits_+x^{\mathbf k} =wt\mathop{\triangleright}\nolimits_c x^{\mathbf k}. \!\!\!\tag*{\qed}
\end{gather*}
\renewcommand{\qed}{}
\end{proof}

Taking $c=1$ in Lemma~\ref{lem:J_c}(b), we settle Proposition~\ref{prop:automorphism}(d).
Proposition~\ref{prop:automorphism} is proved.


\subsection*{Proof of Theorem~\ref{thm:main}} We retain the notation from the proof of
Proposition~\ref{prop:automorphism}.
Let $G=G(m,p,n)$ as in the theorem, and denote $T:=G\cap (\mathbb C^\times)^n$.
Let $e_T=\sum\limits_{t\in T}t\in \mathbb C T\subset \mathbb C G$.
Clearly, $\g ie_T=\g 1e_T$ for all $i=1,\ldots,n$.
Hence $Q_{ij}^{(c)}e_T=\g1e_T$ for all $c\in\mathbb C^\times$, so that $Q_w^{(c)}e_T=t_1^{(\det w)}e_T$ for all $w\in
\mathbb S_n$.
Since, as sets, $G=\{wt \,|\, w\in \mathbb S_n, t\in T\}$, one has $e_G=\sum\limits_{w\in\mathbb S_n}w\cdot e_T$ and
\begin{gather*}
J_c(e_G)= \sum\limits_{w\in \mathbb S_n} J_c(w)e_T = \sum\limits_{w\in \mathbb S_n} wQ_w^{(c)}e_T =\sum\limits_{w\in
\mathbb S_n} wt_1^{(\det w)}e_T = e_{\mu(G)},
\end{gather*}
because, as sets, $\mu(G)=\big\{wt_1^{(\det w)} t \,|\, w\in \mathbb S_n, t\in T\big\}$.
Since a~subgroup $G'$ of $\mathbb G_n$ is uniquely determined by $e_{G'}\in\mathbb C \mathbb G_n$, the group $\mu(G)$ is
a~unique subgroup $G'$ of $\mathbb G_n$ such that $J_c(e_G)=e_{G'}$.

Finally, by Lemma~\ref{lem:J_c}(b), $\rho_c(e_G)=\rho_+(e_{\mu(G)})$.
Setting $c=1$ and using injectivity of $\rho_+$, see Remark~\ref{rem:injectivity}, completes the proof of
Theorem~\ref{thm:main}.


\subsection*{Proof of Theorem~\ref{thm:main2}} Let $G=G(m,p,n)$.
Consider the restriction of $J_\mathbf{i}$, $\mathbf{i}=\sqrt{-1}$ to $R G=\mathbb S_n\cdot R T$ where $T=G\cap (\mathbb
C^\times)^n=\mu(G)\cap (\mathbb C^\times)^n$.
Observe that $J_\mathbf{i}(T)=T$ and
\begin{gather*}
J_\mathbf{i}(s_i)=s_iQ_{s_i}^{(\mathbf{i})}=s_i Q_{i,i+1}^{(\mathbf{i})}
=\frac12 s_i\left((1+\mathbf{i})t_i^{(-1)}+(1-\mathbf{i})t_{i+1}^{(-1)}\right)
=\frac{1+\mathbf{i}}{2}\sigma_i+\frac{1-\mathbf{i}}{2}\sigma_i^{-1},
\end{gather*}
where $\sigma_i=s_i t_i^{(-1)}\in \mu(G)$.
Thus, $J_\mathbf{i}(R G)\subseteq R\mu(G)$, so that the automorphism $J_\mathbf{i}$ of $\mathbb C \mathbb G_n$ restricts
to an isomorphism $R G \xrightarrow{\sim} R \mu(G)$.
Theorem~\ref{thm:main2} is proved.
